\section{Conclusion et perspectives}

Ce travail a permis de mettre en place une plateforme d’observabilité APM cohérente, couvrant l’ingestion, le stockage, l’analyse et la visualisation des données. L’usage de TimescaleDB pour les séries temporelles, associé à des vues continues pour les KPI, apporte un socle performant et maintenable. La supervision couvre les trois piliers de l’observabilité --- métriques (Prometheus), journaux (Loki) et traces distribuées (OpenTelemetry + Tempo) --- consolidés dans Grafana, tandis que l’intégration des sauvegardes pgBackRest avec MinIO améliore la résilience globale. Les jeux de tests Postman/Newman et les scripts de déploiement documentés rendent la solution reproductible.

Sur le plan d’architecture, nous avons validé une configuration en mode cluster adaptée aux besoins de supervision et de charge. Ce choix facilite l’évolutivité, la séparation des responsabilités et la montée en puissance future. Dans un contexte de production, le cloud reste souvent le meilleur compromis en termes de coût, d’élasticité et d’exploitation, tandis que l’on-premise répond surtout aux exigences fortes de souveraineté et de conformité.

\subsection*{Apports d'industrialisation}
Au-delà du socle initial, la plateforme a été industrialisée sur plusieurs axes :
\begin{itemize}
    \item alerting orienté SLO avec règles de \emph{burn-rate} multi-fenêtres sur le budget d’erreur ;
    \item traçage distribué OpenTelemetry exporté vers Tempo, et journaux structurés corrélés dans Loki ;
    \item cycle de vie des données : compression et rétention TimescaleDB, plus un contrôle automatisé de la qualité des données ;
    \item documentation d’API auto-générée (OpenAPI / Swagger) et tests de charge k6 servant de garde-fou de performance ;
    \item chaîne CI/CD sécurisée : lint, tests, scan de vulnérabilités (Trivy), analyse statique (CodeQL), SBOM et signature d’image (Cosign) ;
    \item déploiement Kubernetes via un chart Helm et une livraison GitOps avec Argo CD.
\end{itemize}

\subsection*{Perspectives}
Plusieurs axes peuvent encore enrichir la plateforme :
\begin{itemize}
    \item échantillonnage et quotas de traces pour maîtriser le volume en forte charge ;
    \item autoscaling piloté par les métriques applicatives (par exemple via KEDA) ;
    \item externalisation des secrets (Vault, External Secrets) et rotation automatisée ;
    \item promotion multi-environnements en GitOps (dev / staging / prod) ;
    \item détection d’anomalies sur les séries temporelles et exploitation des \emph{embeddings} pgvector ;
    \item tests de résilience (chaos engineering) sur la bascule primaire/réplica et les pannes de stockage.
\end{itemize}
